\MFQLead{状态估计与变点答案}
\begin{enumerate}[leftmargin=*]
  \item 预测为 $p(s_t\mid y_{1:t-1})$，滤波为 $p(s_t\mid y_{1:t})$，平滑为 $p(s_t\mid y_{1:T})$。只有前两者可按相应时点在线获得。
  \item 交换两套状态参数与转移矩阵行列不改变观测似然；经济命名需要额外排序约束，且应报告该约束。
  \item $K=P^-H/(H^2P^-+R)$；$m^+=m^-+K(y-Hm^-)$；$P^+=(1-KH)P^-$。
  \item $\pi_{t|t-1}=\pi_{t-1|t-1}P$；再计算 $\tilde\pi_{t,j}=f_j(y_t)\pi_{t|t-1,j}$ 并除以总和。
  \item 较大 $Q/R$ 提高增益、缩短响应时间，但也放大观测冲击；敏感性报告应同时给误差与平滑度。
  \item 至少保存每次初值、收敛状态、对数似然、转移矩阵和对齐后的概率路径。近似同似然但不同参数说明弱识别。
  \item 负例把估计截止设为决策日之后，调用时间验证器必须非零退出；修复是只保存当时滤波状态并从下一滚动窗重新估计。
\end{enumerate}

\section*{逐题推导与核对}
\begin{enumerate}[leftmargin=*]
  \item 预测分布为 $p(s_t\mid y_{1:t-1})$，滤波在看到 $y_t$ 后更新为 $p(s_t\mid y_{1:t})$，平滑再使用全样本得到 $p(s_t\mid y_{1:T})$。实时交易只能使用前两者；把平滑概率回填历史就是未来信息。
  \item 对两状态模型交换标签映射 $1\leftrightarrow2$，同时交换初始概率、转移矩阵的行列和发射参数，所有观测似然不变。这是标签交换不识别；若要叫“高波动状态”，必须事先规定按方差排序并在每个滚动窗对齐。
  \item 预测方差为 $P^-=P+Q$，观测残差方差为 $S=H^2P^-+R$。最小化后验均方误差或对联合高斯完成平方，得到 $K=P^-H/S$；更新均值 $m^+=m^-+K(y-Hm^-)$，方差 $P^+=(1-KH)P^-$。代入 $H=1,P^-=1,R=.25,y=1.2$ 得 $K=.8,m^+=.96,P^+=.2$。
  \item 先验状态概率为 $\pi_{t|t-1,j}=\sum_i\pi_{t-1|t-1,i}P_{ij}$；观测密度为 $f_j(y_t)$，未归一化权重 $\tilde\pi_{t,j}=f_j(y_t)\pi_{t|t-1,j}$，归一化为 $\pi_{t|t,j}=\tilde\pi_{t,j}/\sum_k\tilde\pi_{t,k}$。每一步都应记录分母，分母为零即模型/数据失败。
  \item $Q/R$ 大表示状态噪声相对观测噪声大，$K$ 增大、滤波更快跟随新观测但更抖；$Q/R$ 小则路径更平滑、响应更慢。扫描一组比值时同时报告滞后、残差、状态方差和样本外损失，不能只看一条最贴合路径。
  \item 多初值拟合每次保存初值、收敛状态、迭代次数、对数似然、参数和状态路径。先按状态方差/均值对齐标签，再比较似然和路径；近似同似然但参数差异大说明弱识别，应报告不确定性而不是挑最漂亮的一次。
  \item 负例把估计截止设在决策日之后，时间验证器应检查 \texttt{fit\_end <= decision\_time} 并非零退出。修复是只保存当时滤波状态，下一滚动窗重新估计；平滑结果可用于事后诊断但不能进入历史交易。
\end{enumerate}
